\section{Results}
\label{sec:results}

% Nine floats in a few pages.  LaTeX's default float fractions reserve so much
% of each page for text that the tables would be pushed several pages past the
% prose that discusses them; the values below let a page carry more table and
% fewer orphaned strips of text.  They are restored at the end of the section.
\renewcommand{\topfraction}{0.92}
\renewcommand{\bottomfraction}{0.75}
\renewcommand{\textfraction}{0.06}
\renewcommand{\floatpagefraction}{0.80}
\setcounter{topnumber}{3}
\setcounter{bottomnumber}{2}
\setcounter{totalnumber}{4}
\setlength{\textfloatsep}{12pt plus 2pt minus 4pt}
\setlength{\floatsep}{10pt plus 2pt minus 3pt}
\setlength{\intextsep}{12pt plus 2pt minus 3pt}

The result of this work is the removal of the failure mode of
\S\ref{sec:diag-baseline}, and it is not a strength jump. Over the
\numVfiveStyleN{} opponents of Table~\ref{tab:perstyle} the two policies' mean
win rates are \numVfiveMean\% and \numVfourMean\%, their worst cases are
\numVfiveWorst\% and \numVfourWorst\%, and \vfour{} is marginally the better of
the two on minimax regret over that style set. The evidence for \vfive{} is
Table~\ref{tab:pathology}, the forced-endgame measurement of
\S\ref{sec:results-forced}, and the gain of $\numDecWithholderDelta$ points
against the deception archetype that motivated the work, replicated at both of
the banks in Table~\ref{tab:deception}. The same table records a smaller and
equally replicated \emph{loss} of $\numDecFeintDelta$ points against a different
deceptive archetype, and the evidence against a strength claim is that loss
together with every row of Table~\ref{tab:h2h} and of Table~\ref{tab:perstyle}.
All of it is reported below at the same weight.

\subsection{The pathology, on the same instrument}
\label{sec:results-pathology}

\begin{table}[t]
\centering
\caption{\texttt{fish pathology} run on the \vfour{} mirror, the \vfive{} mirror
and the cross match. Population: \numPathologyDeals{} deals per column at both
team orientations; seed \numPathologySeed{} for the two mirrors and
\numPathologyCrossSeed{} for the cross match, full rules dialect, both policies
frozen. The cross match is \numCrossGames{} distinct games. The two mirror
columns are \numFourGames{}: in a mirror the two orientations run identical
policies and play byte-identical games, so the instrument's raw counts are each
game counted twice, and every count in those two columns is reported here per
distinct game --- the same halving C3 in \S\ref{sec:corrections} applies to the
forced-endgame figure. Re-running each bank at one orientation returns exactly
half of every count and every rate to the printed digit, so no rate in this
table is affected. The instrument also pools \emph{both} teams, which is why the
cross-match column is a property of the pair and not of either policy: the dead
asks it counts are \vfour{}'s, and \S\ref{sec:results-forced} returns to the
consequence of that pooling for the forced-endgame row. No intervals; these are
counts and rates over one bank per column. Source:
\texttt{research/v05/results/E2-pathology.txt}.}
\label{tab:pathology}
\footnotesize
\begin{tabular}{lrrr}
\toprule
& \textbf{\vfour{} mirror} & \textbf{\vfive{} mirror} & \textbf{cross} \\
\midrule
events per game                        & \numFourEventsMean{}  & \numFiveEventsMean{}  & \numCrossEventsMean{} \\
\quad median / p90 / p99               & \numFourEventsMedian{}\,/\,\numFourEventsPninety{}\,/\,\numFourEventsPninetynine{}
                                       & \numFiveEventsMedian{}\,/\,\numFiveEventsPninety{}\,/\,\numFiveEventsPninetynine{}
                                       & \numCrossEventsMedian{}\,/\,\numCrossEventsPninety{}\,/\,\numCrossEventsPninetynine{} \\
asks per game                          & \numFourAsks{}        & \numFiveAsks{}        & \numCrossAsks{} \\
ask hit rate                           & \numFourHit\%         & \numFiveHit\%         & \numCrossHit\% \\
\midrule
provably dead asks                     & \numFourDeadAsk\%     & \numFiveDeadAsk\%     & \numCrossDeadAsk\% \\
dead runs                              & \numFourDeadRuns{}    & \numFiveDeadRuns{}    & \numCrossDeadRuns{} \\
\quad longest                          & \numFourDeadRunLongest{} & \numFiveDeadRunLongest{} & \numCrossDeadRunLongest{} \\
games with a dead run $\geq 6$         & \numFourDeadRunGames\% & \numFiveDeadRunGames\% & \numCrossDeadRunGames\% \\
starved turns                          & \numFourStarved\%     & \numFiveStarved\%     & \numCrossStarved\% \\
exact repeat (actor, card, target)     & \numFourRepeatAsk\%   & \numFiveRepeatAsk\%   & \numCrossRepeatAsk\% \\
asks inside a half-suit the team owns & \numFourOwnLocked\%   & \numFiveOwnLocked\%   & \numCrossOwnLocked\% \\
\midrule
declarations                           & \numFourDeclN{}       & \numFiveDeclN{}       & \numCrossDeclN{} \\
\quad wrong                            & \numFourDeclWrong\%   & \numFiveDeclWrong\%   & \numCrossDeclWrong\% \\
at or after the forcing horizon        & \numFourPostHorizonDecl{} & \numFivePostHorizonDecl{} & \numCrossPostHorizonDecl{} \\
\quad of those, wrong                  & \numFourPostHorizonWrong\% & --- & --- \\
forced-endgame declarations (pooled)   & \numFourForced{}      & \numFiveForced{}      & \numCrossForced{} \\
\quad of those, wrong                  & \numFourForcedWrong\% & \numFiveForcedWrong\% & \numCrossForcedWrong\% \\
\midrule
games reaching the action cap          & \numFourCapRate\%     & \numFiveCapRate\%     & \numCrossCapRate\% \\
\bottomrule
\end{tabular}
\end{table}

Table~\ref{tab:pathology} is the paper's result. The three statistics promoted to
commit gates in \S\ref{sec:mech-built} --- longest dead run, share of games
containing a dead run of six, and declarations made at or after the forcing
horizon --- read \numFiveDeadRunLongest{}, \numFiveDeadRunGames\% and
\numFivePostHorizonDecl{} on the \vfive{} mirror bank, against
\numFourDeadRunLongest{}, \numFourDeadRunGames\% and \numFourPostHorizonDecl{}
for \vfour{} on the identical bank. The provably-dead ask rate falls from
\numFourDeadAsk\% to \numFiveDeadAsk\%, which is a property of the gate rather
than a measurement: the predicate of Definition~\ref{def:dead} is the same
deduction the filter runs, so a gated policy can record a dead ask only through
the starved-turn fallback, and on this bank the fallback never fired at all
(\numFiveStarved\% of asks).

What is a measurement, and what the gate did not guarantee, is the rest of the
table. Ask hit rate rises from \numFourHit\% to \numFiveHit\%, and ask volume
falls from \numFourAsks{} to \numFiveAsks{} per game: removing a class of moves
that could not succeed did not merely shorten the games, it changed what the
surviving asks are spent on. The event distribution moves at the tail and not
only at the mean, which is the property the swept holding cost of
\S\ref{sec:mech-rejected} failed to obtain --- the 99th percentile falls from
\numFourEventsPninetynine{} events to \numFiveEventsPninetynine{}, and the median
from \numFourEventsMedian{} to \numFiveEventsMedian{}. Declaration error falls
from \numFourDeclWrong\% to \numFiveDeclWrong\%, and the \numFourPostHorizonDecl{}
post-horizon declarations that were \numFourPostHorizonWrong\% wrong have no
counterpart because no game reaches the horizon.

Three rows are worth reading against the natural expectation rather than for it.
Exact repeats fall to \numFiveRepeatAsk\% but not to zero, and should not: cards
move, so re-asking a (card, target) pair can be the best move on the board once a
public transfer has put the card where it was not before. That is the same
observation that sank the repetition guard (\S\ref{sec:results-ablations}). Asks
inside a half-suit the actor's own team already owns fall from \numFourOwnLocked\%
to \numFiveOwnLocked\%, and again not to zero: those asks are legal, are not
provably dead when the actor cannot establish the ownership, and emit ask-legality
certificates. \vfive{} does not choose them for that reason --- the net-information
term M3 is unbuilt --- so the remaining rate is incidental, exactly as
\S\ref{sec:limitations} says of the \vfour{} rate. Finally, the cross-match column
records \numCrossDeadAsk\% dead asks with every dead run of length
\numCrossDeadRunLongest{}: the instrument pools both teams, these asks are
\vfour{}'s, and the run lengths are one because the seat that would have continued
the cycle is now gated.

\subsection{Engine and information-safety verification}
\label{sec:results-verify}

\texttt{fish verify} plays the ordered cross-product of the engine's
\numVerifyPool{}-policy pool --- \numVerifyPairs{} ordered pairs over a budget of
\numVerifyGames{} games --- with the knowledge audit active, and records
\numVerifyViolations{} violations in \numVerifyChecks{} soundness checks: no
derived knowledge state, capacity or certificate was ever contradicted by the
actual deal. Card conservation holds in every game, with \numVerifySetFail{}
set-conservation failures; \numVerifyLimitGames{} games reach the action limit;
and bit-identical replay is \numVerifyDeterminism{}.

The scope of that sweep needs stating, because it is narrower than it looks. The
pool it plays is \vfour{}, \vthree{}, \vtwo{} and the six scripted styles, and it
does \emph{not} contain \vfive{}; the determinism check is run on a
\vfour{}-against-\vthree{} bank. What E1 establishes is that the rules driver, the
public deduction state and the certificate machinery are sound --- which is the
machinery M1's gate is a pure function of, so a sound deduction state is exactly
the precondition the gate needs --- and not that \vfive{}'s own play was audited.
For \vfive{} the corresponding evidence is the action-limit rate, which is
\numFiveCapRate\% for the \vfive{} mirror and \numFourCapRate\% for \vfour{} on
the bank of Table~\ref{tab:pathology}, and is non-zero on \numLimitNonzero{} of
the \numLimitRows{} match rows behind the tables in this section: no result here
is an artefact of the cap ending a game. Sources:
\texttt{research/v05/results/E1-verify.txt}, \texttt{engine/src/main.cpp}.

\subsection{The forced endgame, at volume}
\label{sec:results-forced}

\begin{table}[t]
\centering
\caption{Forced-endgame declarations, per declaring team (E8). Population:
\numForcedEightDeals{} deals in six seat rotations, \numForcedEightGames{} games
per arm, seed \numForcedEightSeed{}, \vfive{} against \vfour{}, both frozen. The
two arms are the two sides of the same matches, scored separately; ``rate'' is
forced declarations per game by that team and ``correct'' is the share of them
that named the true allocation. Declaration totals are all declarations by that
team, forced and voluntary. No intervals: a forced declaration is a rare event
and these are pooled counts. Source:
\texttt{research/v05/results/E8-forced-endgame.txt}.}
\label{tab:forced}
\small
\begin{tabular}{lrrrr}
\toprule
& \multicolumn{2}{c}{\textbf{forced declarations}} & \multicolumn{2}{c}{\textbf{all declarations}} \\
\cmidrule(lr){2-3}\cmidrule(lr){4-5}
& \textbf{rate/game} & \textbf{correct} & \textbf{rate/game} & \textbf{correct} \\
\midrule
\vfour{} & \numVfourForcedRate{} & \numVfourForcedAcc\% & \numVfourDeclRateEight{} & \numVfourDeclAccEight\% \\
\vfive{} & \numVfiveForcedRate{} & \numVfiveForcedAcc\% & \numVfiveDeclRateEight{} & \numVfiveDeclAccEight\% \\
\bottomrule
\end{tabular}
\end{table}

The forced endgame is too rare for the pathology bank of
Table~\ref{tab:pathology} to resolve. Its \vfive{} mirror column counts
\numFiveForced{} of these declarations in \numFiveGames{} games, which is a
denominator that supports no rate. E8 therefore plays
\numForcedEightGames{} games per arm and scores each team separately.
Table~\ref{tab:forced} gives the result. \vfive{} enters the forced endgame
\numForcedRateRatio$\times$ less often --- \numVfiveForcedRate{} declarations per
game against \numVfourForcedRate{} --- and when it does, it names a
capacity-feasible allocation that is correct \numVfiveForcedAcc\% of the time
(\numVfiveForcedRightN{} of \numVfiveForcedN{} forced declarations), against
\numVfourForcedAcc\% for \vfour{} (\numVfourForcedRightN{} of
\numVfourForcedN{}). The denominators are worth carrying: the \vfive{} rate rests
on \numVfiveForcedN{} events, which is enough to separate it from \vfour{}'s and
not enough to place it to the nearest point. E8 does not decompose the two
effects by mechanism, so the attribution that follows is the design's and not
this experiment's: M1 is what removes the extended runs of guaranteed misses
during which a team is stripped of cards while half-suits remain undeclared,
which is how the cardless endgame was being reached, and M2 is what stops the
declarer, once there, naming an allocation that gives a teammate more cards than
that teammate can hold.

\numVfiveForcedAcc\% is not the ceiling. \S\ref{sec:diag-forced} measures the best
capacity-feasible allocation to be correct \numFeasibleRight\% of the time under
the reference posterior. That ceiling was measured on \vfour{}'s mirror
forced-endgame states, not on the states \vfive{} reaches here, so it bounds this
number only to the extent that the two populations resemble each other; read as
an indication rather than a bound, roughly two-fifths of the available accuracy
remains unclaimed. The gap is a live target rather than a residual: M2 ranks
feasible assignments by an estimate computed on the deployed approximate belief,
and it does not attempt the per-seat knowledge model (M4) that would sharpen the
ranking.

\paragraph{Why this is reported per declaring team.} \texttt{fish pathology}
pools both teams' forced declarations into a single count. In the cross match of
Table~\ref{tab:pathology} that pooled figure is \numCrossForced{} declarations,
\numCrossForcedWrong\% wrong, and it is a property of neither policy:
Table~\ref{tab:forced} scores the same kind of event separately by declaring team,
on a far larger bank, and finds \numVfiveForcedAcc\% correct for
\vfive{} against \numVfourForcedAcc\% for \vfour{}, so the pooled cross-match
figure is dominated by \vfour{}'s side of those games. This is the same class of
error corrected at C3 in \S\ref{sec:corrections}, where a mirror run's six seat
rotations were read as six independent samples although they are three deal
shifts times two team orientations and both orientations play byte-identical
games. Forced-endgame accuracy is attributable only when it is reported per
declaring team, which is why Table~\ref{tab:forced}, and not the pathology row,
is the measurement of record.

\subsection{Head-to-head against \vfour{}, on deals never fitted against}
\label{sec:results-h2h}

\begin{table}[t]
\centering
\caption{\vfive{} against \vfour{}, five independent evaluation banks (E3).
Population: \numHeadBankDeals{} deals in six seat rotations per bank,
\numHeadBankGames{} games per row and \numHeadGames{} in total, full rules dialect,
both policies frozen. Intervals are deal-clustered percentile bootstrap. The
deals are held out from the fit of \S\ref{sec:fitting}; \vfour{} itself was a
member of the fitting panel, so the opponent identity is not. Source:
\texttt{research/v05/results/E3-headtohead.jsonl}.}
\label{tab:h2h}
\small
\begin{tabular}{lrr}
\toprule
\textbf{Bank seed} & \textbf{\vfive{} win rate} & \textbf{95\% CI} \\
\midrule
\numHeadSeedOne{}   & \numHeadWinOne\%   & \numHeadCIOne\% \\
\numHeadSeedTwo{}   & \numHeadWinTwo\%   & \numHeadCITwo\% \\
\numHeadSeedThree{} & \numHeadWinThree\% & \numHeadCIThree\% \\
\numHeadSeedFour{}  & \numHeadWinFour\%  & \numHeadCIFour\% \\
\numHeadSeedFive{}  & \numHeadWinFive\%  & \numHeadCIFive\% \\
\midrule
mean of the five    & \numHeadWin\%      & range \numHeadRange\% \\
\bottomrule
\end{tabular}
\end{table}

Across the five banks of Table~\ref{tab:h2h}, \vfive{} wins \numHeadWin\% of
\numHeadGames{} games against \vfour{}, with individual banks spanning
\numHeadRange\% and every interval containing $50\%$. The mean margin is about a
point, which is smaller than the spread between banks, and no claim of a strength
advantage is made from it.

These are not the only banks in this report on which the frozen configuration
meets \vfour{}. Counting the per-style cell of Table~\ref{tab:perstyle}
(\numPanelVfiveVfour\%, the figure the abstract and \S\ref{sec:intro} quote
because it is the cell the worst case of that table is read from), the shipped
row of the ablation bank (\numAblFull\%) and the default dialect of
Table~\ref{tab:dialects-v05} (\numDialectDefault\%), \numHeadBanks{} independent
banks measure the same pairing and they span \numHeadBankSpan\%. That span, and
not any one of its members, is what this report knows about the head-to-head:
it straddles even money, its width exceeds every margin inside it, and its
lowest member, \numHeadBankLow\%, puts \vfive{} below \vfour{}.

\paragraph{The fitting regression, stated plainly.} The generation the fit of
\S\ref{sec:fitting} selected from scores \numFitWorst\% against \vfour{}
in-panel, on the banks it was being optimised on. Held out, after freezing, the
same configuration scores \numHeadWin\%. In-panel \numFitWorst\% versus held-out
\numHeadWin\% is the regression, and the held-out figure is the one reported
here and everywhere else in this paper. Nothing about the in-panel number should
be quoted as a result.

The composition of the margin is worth recording because it is not the
composition one would predict from the pathology table. \vfive{} takes
\numHeadSetsFive{} half-suits per game against \numHeadSetsFour{}, and declares
slightly more often --- \numHeadDeclRateFive{} declarations per game against
\numHeadDeclRateFour{} --- at a slightly \emph{lower} accuracy,
\numHeadDeclAccFive\% against \numHeadDeclAccFour\%. \vfour{} gets
\numHeadDeclErrFour\% of its declarations wrong in these games, against
\numFourDeclWrong\% in its own mirror, because against a \vfive{} opponent the
games end at \numHeadEvents{} events and \vfour{} never reaches the horizon that
forces its bad declarations.
The failure mode of \S\ref{sec:diag-baseline} is a property of the interaction,
not of one seat's rule in isolation, and repairing one side of the match removes
most of it from both.

\subsection{The per-style profile, and the worst case}
\label{sec:results-style}

\begin{table}[t]
\centering
\caption{Per-opponent win rate for both policies against the same
\numVfiveStyleN{}-member style set (E4). Population: \numStyleDeals{} deals in six
seat rotations, \numStyleGames{} games per cell, seed \numStyleSeed{}, full rules
dialect, both policies frozen. $\Delta$ is \vfive{} minus \vfour{} in percentage
points. The \vfour{}-against-\vfour{} cell is $50\%$ by construction, not by
measurement. Minimax regret is computed against the better of the two
\emph{measured} arms on each style, not against a per-style optimum, which was
not computed (\S\ref{sec:limitations}). Source:
\texttt{research/v05/results/E4-perstyle.jsonl}.}
\label{tab:perstyle}
\small
\begin{tabular}{lrrr}
\toprule
\textbf{Opponent} & \textbf{\vfive{}} & \textbf{\vfour{}} & \textbf{$\Delta$} \\
\midrule
\vfour{}      & \numPanelVfiveVfour\%       & \numPanelVfourVfour\%       & $\numDeltaVfour$ \\
\vthree{}     & \numPanelVfiveVthree\%      & \numPanelVfourVthree\%      & $\numDeltaVthree$ \\
\vtwo{}       & \numPanelVfiveVtwo\%        & \numPanelVfourVtwo\%        & $\numDeltaVtwo$ \\
lockout       & \numPanelVfiveLockout\%     & \numPanelVfourLockout\%     & $\numDeltaLockout$ \\
detective     & \numPanelVfiveDetective\%   & \numPanelVfourDetective\%   & $\numDeltaDetective$ \\
diversifier   & \numPanelVfiveDiversifier\% & \numPanelVfourDiversifier\% & $\numDeltaDiversifier$ \\
hunter        & \numPanelVfiveHunter\%      & \numPanelVfourHunter\%      & $\numDeltaHunter$ \\
bluffer       & \numPanelVfiveBluffer\%     & \numPanelVfourBluffer\%     & $\numDeltaBluffer$ \\
random        & \numPanelVfiveRandom\%      & \numPanelVfourRandom\%      & $\numDeltaRandom$ \\
\midrule
\textbf{worst case}          & \textbf{\numVfiveWorst\%} & \textbf{\numVfourWorst\%} & \\
\quad attained on            & \numVfiveWorstOpp{}       & \numVfourWorstOpp{}       & \\
mean over the style set      & \numVfiveMean\%           & \numVfourMean\%           & \\
minimax regret               & \numVfiveRegret{}         & \numVfourRegret{}         & \\
\quad attained on            & \numVfiveRegretOpp{}      & \numVfourRegretOpp{}      & \\
\bottomrule
\end{tabular}
\end{table}

Table~\ref{tab:perstyle} is the table this project's standing evaluation
preference asks for, and it says that on win rate the two policies are level.
\vfive{} is ahead on four styles, behind on four, and identical on one; the
largest gain is $\numDeltaLockout$ against the turn-starvation lockout and the
largest loss $\numDeltaVtwo$ against \vtwo{}. The two means differ by
\numStyleMeanGap{} points, computed from the unrounded cells. On minimax regret \vfour{} is marginally the
better policy, \numVfourRegret{} against \numVfiveRegret{}, its regret attained on
\numVfourRegretOpp{} and \vfive{}'s on \numVfiveRegretOpp{}. Neither figure
supports an ordering.

\paragraph{The worst case, and what \vfour{}'s published one excluded.} The v0.4
study reported that \vfour{}'s lowest win rate against any panel member was
\numVfourVsVthree\%. That panel contained no opponent of \vfour{}'s own strength.
Once one is included, the worst case for both policies falls to about $50\%$:
\numVfiveWorst\% for \vfive{} and \numVfourWorst\% for \vfour{}, both attained on
\numVfiveWorstOpp{}. This is not a defect that a better policy in this family
would repair. A policy's mirror match is a two-team game between identical
strategies, so \vfour{}'s own cell is exactly $50\%$ by construction, and any
policy strong enough that no member of a fixed panel beats it will have its worst
case set by a copy of itself. A worst-case figure quoted over a panel that
excludes the mirror is therefore a statement about the panel. C9 in
\S\ref{sec:corrections} records the scope correction; the operational consequence
is M10's, and it is that the mirror belongs in the panel that the worst case is
computed over as well as in the panel that is fitted against.

\subsection{The deception panel}
\label{sec:results-deception}

\begin{table}[t]
\centering
\caption{\vfive{} and \vfour{} against three deceptive archetypes (E10), at two
independent seed banks. Population: \numDecDeals{} deals in six seat rotations,
\numDecGames{} games per cell, seeds \numDecSeedA{} and \numDecSeedB{}, full
rules dialect, both policies frozen. All three are commitment deceivers: they
follow a fixed deceptive rule for the whole game. \emph{Silent} never asks in the
half-suit it holds most cards of until forced --- the uncapped variant, not the
hit-probability-capped one whose figure \S\ref{sec:diag-prior} quotes;
\emph{feint} preferentially asks in a half-suit it holds exactly one card of,
manufacturing a misleading ask-legality certificate when the material cost is
small; \emph{withholder} is the project owner's own manoeuvre, declining to ask in
a half-suit for several turns after being asked in it while holding other cards of
it. $\Delta$ is the mean over both banks of \vfive{} minus \vfour{}. Source:
\texttt{research/v05/results/E10-deception.md}.}
\label{tab:deception}
\small
\begin{tabular}{lrrrrr}
\toprule
& \multicolumn{2}{c}{\textbf{seed \numDecSeedA{}}} & \multicolumn{2}{c}{\textbf{seed \numDecSeedB{}}} & \\
\cmidrule(lr){2-3}\cmidrule(lr){4-5}
\textbf{Opponent} & \textbf{\vfive{}} & \textbf{\vfour{}} & \textbf{\vfive{}} & \textbf{\vfour{}} & \textbf{$\Delta$} \\
\midrule
silent      & \numDecSilentFiveA\%      & \numDecSilentFourA\%      & \numDecSilentFiveB\%      & \numDecSilentFourB\%      & $\numDecSilentDelta$ \\
feint       & \numDecFeintFiveA\%       & \numDecFeintFourA\%       & \numDecFeintFiveB\%       & \numDecFeintFourB\%       & $\numDecFeintDelta$ \\
withholder  & \numDecWithholderFiveA\%  & \numDecWithholderFourA\%  & \numDecWithholderFiveB\%  & \numDecWithholderFourB\%  & $\numDecWithholderDelta$ \\
\bottomrule
\end{tabular}
\end{table}

The withholder is a direct model of the manoeuvre reported from live play that
started this work: hold cards of a half-suit you have been asked for, then
decline to ask back in it. Against that archetype \vfive{} gains
$\numDecWithholderDelta$ points, and the two banks agree closely on the size of
it --- $\numDecWithholderDeltaA$ and $\numDecWithholderDeltaB$ --- which is the
strongest replication anywhere in this report. It also gains
$\numDecSilentDelta$ against the silent archetype, but there only the sign
replicates: the banks read $\numDecSilentDeltaA$ and $\numDecSilentDeltaB$, and
the first of those is a fraction of the width of a bank, so the silent result
should be read as a direction and not as a magnitude. The mechanism is not a better
opponent model --- M3 through M7 are unbuilt and the two policy-prior scalars are
the same construct in both arms. It is that \vfive{} no longer spends turns on
asks it can prove will miss, so a misleading \emph{absence} of asks has much less
leverage over the position. The manoeuvre works by making the deceiver's holdings
appear to lie elsewhere, and a policy already spending \numFourDeadAsk\% of its
asks on cards it could prove were unavailable was donating turns for that
misdirection to work in.

\paragraph{The loss, and its cause.} Against the feint \vfive{} is
$\numDecFeintDelta$ points worse, and that replicates at both banks too
($\numDecFeintDeltaA$ and $\numDecFeintDeltaB$). The
feint manufactures a false-positive ask-legality certificate rather than
withholding a true one, and the fit raised \texttt{priorTheta} --- the weight on
``this player asked in this half-suit'' --- from \numPriorThetaFour{} to
\numPriorThetaFive{}. \vfive{} therefore weights a manufactured certificate more
heavily than \vfour{} did. This is the exposure the diagnosis predicted:
\S\ref{sec:diag-prior} establishes that the vulnerability is over-weighting the
policy prior, not weighting it at all, and the corresponding removal is worse
still --- deleting both prior scalars costs \numPriorDeleteCost{} points
[\numPriorDeleteCI] on a panel containing no honest opponent, against the
\numAgnosticCost{} points \S\ref{sec:mech-rejected} records on the mixed panel.

\paragraph{The parameter is not identified, and was not tuned.} A sweep of
\texttt{priorTheta} over \numThetaSweepN{} values, $\{\numThetaSweep\}$, at
\numDecSeeds{} seeds does not identify it: the ordering of the settings is not
stable between the two banks and every pairwise difference sits inside the
cluster-bootstrap interval. The fitted value was therefore retained rather than
adjusted, because adjusting it on this evidence would be selection on noise. The
principled repair is M7 --- a per-seat online type posterior with data-biased
shrinkage, so that a certificate from a seat whose type posterior has drifted
toward ``deceptive'' is discounted --- and M7 is specified and not built
(\S\ref{sec:mech-designed}).

\paragraph{Consequence for the worst case.} Across the full
\numTwelveStyles{}-style set --- the \numVfiveStyleN{} of
Table~\ref{tab:perstyle} plus these three --- the lowest cell \vfive{} records
anywhere is \numDecWorst\% against the \numDecWorstOpp{}, on the weaker of the
two deception banks, marginally below its mirror cell of \numVfiveWorst\%.
Which of the two is \emph{the} worst case is not resolved by these banks:
averaged over both deception banks the \numDecWorstOpp{} reads
\numDecFeintFive\%, which is above the mirror rather than below it, and the two
candidates sit within half a point of each other either way. The honest
statement is that \vfive{}'s floor over the twelve styles is about
\numVfiveWorst\%, attained on the mirror or on the \numDecWorstOpp{} depending on
the bank. \vfour{}'s worst case remains its own mirror at \numVfourWorst\%.
Neither policy falls below a coin flip against anything in the set, and no
aggregate over it should be quoted without this table and
Table~\ref{tab:perstyle} beside it.

\subsection{Mechanism ablations}
\label{sec:results-ablations}

\begin{table}[t]
\centering
\caption{Frozen-parameter mechanism ablations (E5). Population:
\numAblFiveDeals{} deals in six seat rotations, \numAblFiveGames{} games per row,
seed \numAblFiveSeed{}, every row played against frozen \vfour{}, full rules
dialect. Every row runs \vfive{}'s fitted parameter vector; only the mechanism
switches change, so the control is \emph{not} \vfour{} --- it is the \vfive{}
vector with the mechanisms disabled, which is why it differs from the
\vfour{}-parameter control of Table~\ref{tab:mechanisms}. The last two rows
\emph{add} a mechanism to the shipped configuration rather than removing one.
``Decl.\ acc.'' is the ablated arm's own declaration accuracy and ``opp.'' is
\vfour{}'s in the same games. Intervals are deal-clustered percentile bootstrap.
Source: \texttt{research/v05/results/E5-ablations.jsonl}.}
\label{tab:ablations-v05}
\footnotesize
\begin{tabular}{llrrrrr}
\toprule
\textbf{Switches} & \textbf{Mechanisms} & \textbf{win rate} & \textbf{95\% CI}
  & \textbf{events} & \textbf{decl.\ acc.} & \textbf{opp.} \\
\midrule
\texttt{m1=0,m2=0,stage2=1} & control          & \numAblControl\%    & \numAblControlCI\%    & \numAblControlEvents{}    & \numAblControlDeclAcc\%    & \numAblControlOppDeclAcc\% \\
\texttt{m1=1}               & M1               & \numAblMone\%       & \numAblMoneCI\%       & \numAblMoneEvents{}       & \numAblMoneDeclAcc\%       & \numAblMoneOppDeclAcc\% \\
\texttt{m2=1}               & M2               & \numAblMtwo\%       & \numAblMtwoCI\%       & \numAblMtwoEvents{}       & \numAblMtwoDeclAcc\%       & \numAblMtwoOppDeclAcc\% \\
\texttt{stage2=0}           & M8               & \numAblMeight\%     & \numAblMeightCI\%     & \numAblMeightEvents{}     & \numAblMeightDeclAcc\%     & \numAblMeightOppDeclAcc\% \\
\texttt{m1=1,m2=1}          & M1 + M2          & \numAblMoneMtwo\%   & \numAblMoneMtwoCI\%   & \numAblMoneMtwoEvents{}   & \numAblMoneMtwoDeclAcc\%   & \numAblMoneMtwoOppDeclAcc\% \\
\texttt{m1=1,stage2=0}      & M1 + M8          & \numAblMoneMeight\% & \numAblMoneMeightCI\% & \numAblMoneMeightEvents{} & \numAblMoneMeightDeclAcc\% & \numAblMoneMeightOppDeclAcc\% \\
\texttt{m2=1,stage2=0}      & M2 + M8          & \numAblMtwoMeight\% & \numAblMtwoMeightCI\% & \numAblMtwoMeightEvents{} & \numAblMtwoMeightDeclAcc\% & \numAblMtwoMeightOppDeclAcc\% \\
shipped \vfive{}            & M1 + M2 + M8     & \numAblFull\%       & \numAblFullCI\%       & \numAblFullEvents{}       & \numAblFullDeclAcc\%       & \numAblFullOppDeclAcc\% \\
\midrule
\texttt{+m1p=1}             & \quad + ownership scaled by $p$ & \numAblOwnershipP\%  & \numAblOwnershipPCI\%  & \numAblOwnershipPEvents{}  & \numAblOwnershipPDeclAcc\%  & \numAblOwnershipPOppDeclAcc\% \\
\texttt{+norepeat=1}        & \quad + repetition guard        & \numAblRepeatGuard\% & \numAblRepeatGuardCI\% & \numAblRepeatGuardEvents{} & \numAblRepeatGuardDeclAcc\% & \numAblRepeatGuardOppDeclAcc\% \\
\bottomrule
\end{tabular}
\end{table}

Every row of Table~\ref{tab:ablations-v05} is a frozen-parameter ablation: the
fitted vector is \vfive{}'s in all of them, and no row was refitted around the
mechanism it removes. A cheap removal may still be doing work a refitted policy
would find another way to do, and an expensive removal may be expensive only
because the surrounding weights were fitted to rely on it. That qualification
applies to every row and is not repeated per row; a matched-budget refit is one of
the experiments \S\ref{sec:limitations} names as not run.

\paragraph{Stage 2 is dead code once M1 is on.} The \texttt{m1=1} and
\texttt{m1=1,stage2=0} rows are identical to the digit in win rate, events,
declaration accuracy and their opponent's declaration accuracy, and so are the
\texttt{m1=1,m2=1} and shipped rows. This is not a coincidence of the bootstrap:
with the live-ask gate on, games end at about \numAblFullEvents{} events and the
event-count escalation's second stage is never reached, so switching it off
changes nothing that happens. The measured cost of \numStageTwoCost{} points that
\S\ref{sec:diag-dilemma} attributes to that rule is a cost it can only levy on a
policy that deadlocks.

\paragraph{The rows with the best win rates are the rows that still deadlock.}
M8 alone is the largest number in the table, \numAblMeight\%, and M2 + M8 the
second largest, \numAblMtwoMeight\%; both play games of about
\numAblMeightEvents{} events against the control's \numAblControlEvents{}, so both
retain the failure mode in full. They win by declining to cash a declaration the
policy itself scores near zero: their own declaration accuracy rises from the
control's \numAblControlDeclAcc\% to \numAblMeightDeclAcc\%, while \vfour{}'s in
the same games falls from \numAblControlOppDeclAcc\% to
\numAblMeightOppDeclAcc\% --- an asymmetry this experiment records without
explaining, since the two arms play games of nearly the same length. This is the
forcing dilemma of
\S\ref{sec:diag-dilemma} measured again inside \vfive{}'s switch set: patience
without termination buys points and keeps the pathology, and the configurations
that terminate (\numAblFullEvents{} events, \numAblFullDeclAcc\% own declaration
accuracy, \numAblFullOppDeclAcc\% for the opponent) are the ones that give the
points back. The mechanism package is bought with win rate, not for it.

\paragraph{Two designed mechanisms were rejected on measurement.} Both are
retained as switches, both ship off, and both are recorded here so they are not
rebuilt. Scaling the four ownership features of
\S\ref{sec:mech-built} by hit probability on top of the gate --- \texttt{m1p},
which is a switch and not part of the shipped configuration --- costs $\numAblOwnershipPDelta$ points relative to the
shipped configuration, and was measured at $-\numMonePReject$ points on top of
M1, M2 and M8 at design time, which is why the switch ships off. Its own
diagnostic points at the reason: it \emph{raises} the hit rate, from
\numAblFullHit\% to \numAblOwnershipPHit\%, while losing games. Once the gate has
removed the $p = 0$ case there is no dead ask left for the scaling to suppress,
so all it does is rescale live candidates against one another, and it buys
immediate hits at the expense of the ownership progress the rest of the linear
score was fitted to trade against.

The repetition guard --- \texttt{norepeat}, forbidding a seat from asking the same
(card, target) pair twice --- costs $\numAblRepeatGuardDelta$ points, and was
measured at $-\numGuardReject$ points on top of M1, M2 and M8 at design time. It
also ships off, and it has an explanation worth stating, because the guard looks
like a structural backstop that could not hurt. The reason it does is that cards
move: forbidding a (card, target) repeat forbids asking for a card the actor has
just watched transfer to that opponent in a public event, which is frequently the
best ask on the board. The cost shows in the hit rate, which falls from
\numAblFullHit\% to \numAblRepeatGuardHit\%, and not in the event count, which is
unchanged at \numAblRepeatGuardEvents{} against \numAblFullEvents{}: the guard is
not preventing a deadlock --- there is none left to prevent --- it is deleting
asks the policy would otherwise land. And the
case it was designed to catch is already covered: a repeat of an ask that missed
is provably dead in the sense of Definition~\ref{def:dead}, so M1 removes it
without a separate rule, which is why the guard measured as exactly neutral when
tested alone in Table~\ref{tab:mechanisms} and negative once M1 was present.

\subsection{Calibration}
\label{sec:results-calib}

\begin{table}[t]
\centering
\caption{Aggregate forecast reliability for \vfive{} (E6). Population: the
calibration bank of \numCalibFiveDeals{} deals, seed \numCalibFiveSeed{},
\vfive{} against \vfour{}, weights frozen; the unit is one forecast and $n$
counts forecasts of that type. Voluntary and forced declarations are pooled.
All entries are point estimates: no clustered uncertainty is computed for Brier
score, log loss or expected calibration error, so small differences between them
are not resolved by these numbers. Source:
\texttt{research/v05/results/E6-calibration-v05.txt}.}
\label{tab:calib-v05}
\small
\begin{tabular}{lrrrrrr}
\toprule
\textbf{Forecast} & \textbf{$n$} & \textbf{Brier} & \textbf{log loss}
  & \textbf{ECE} & \textbf{mean pred.} & \textbf{mean obs.} \\
\midrule
ask succeeds        & \numCalibAskN{}  & \numCalibAskBrier{}  & \numCalibAskLogloss{}  & \numCalibAskECE{}  & \numCalibAskPred\%  & \numCalibAskObs\% \\
declaration correct & \numCalibDeclN{} & \numCalibDeclBrier{} & \numCalibDeclLogloss{} & \numCalibDeclECE{} & \numCalibDeclPred\% & \numCalibDeclObs\% \\
\bottomrule
\end{tabular}
\end{table}

Neither aggregate in Table~\ref{tab:calib-v05} characterises the estimator it
scores, for the reason the v0.4 study gave: both are dominated by one bin. Of
\numCalibDeclN{} declaration forecasts, \numCalibDeclTopShare\% fall in the top
decile at \numCalibDeclTopPred\% predicted against \numCalibDeclTopObs\%
realised, which is most of what the expected calibration error of
\numCalibDeclECE{} measures. The decision-relevant bins run the other way: in
$[0.7, 0.8)$ the declaration forecast is \emph{under}confident,
\numCalibDeclMidPred\% predicted against \numCalibDeclMidObs\% realised on
\numCalibDeclMidN{} forecasts: the estimate the stopping rule thresholds is
conservative in exactly the range where the threshold binds.

The ask forecast is the one this report has a reason to look at, because M1
gates on a hard deduction while the ask score prices everything above it. Its
aggregate is close to unbiased --- \numCalibAskPred\% mean predicted against
\numCalibAskObs\% realised over \numCalibAskN{} forecasts --- and it is
overconfident precisely at the bottom, where the lowest decile predicts
\numCalibAskLowPred\% against \numCalibAskLowObs\% realised on
\numCalibAskLowN{} forecasts, and underconfident in the middle, where $[0.4, 0.5)$
predicts \numCalibAskMidPred\% against \numCalibAskMidObs\% on
\numCalibAskMidN{}. The gate itself is unaffected, since it is a deduction and
not a forecast; what the bottom-decile bias means is that among the asks that
survive the gate, the near-dead ones are still priced somewhat too high, which is
a reason to expect the M3 net-information term to have something to price rather
than a defect this report repairs.

\subsection{Rule dialects}
\label{sec:results-dialects}

\begin{table}[t]
\centering
\caption{\vfive{} against \vfour{} under four rules dialects (E7). Population:
\numDialectDeals{} deals in six seat rotations, \numDialectGames{} games per row,
seed \numDialectSeed{}, both policies frozen. ``Forced rate'' is \vfive{}'s
forced declarations per game and the two accuracy columns are the share of each
team's forced declarations that were correct. The \vthree{} legacy dialect
implies declaration on turn only, so its rows track the no-out-of-turn dialect. In
the default and no-cardless dialects a forced declaration occurs a handful of
times per row, so those two accuracy cells rest on single-digit denominators and
should not be read as rates. Intervals are deal-clustered percentile bootstrap.
Source:
\texttt{research/v05/results/E7-rules.jsonl}.}
\label{tab:dialects-v05}
\footnotesize
\begin{tabular}{lrrrrr}
\toprule
& & & \textbf{forced} & \multicolumn{2}{c}{\textbf{forced correct}} \\
\cmidrule(lr){5-6}
\textbf{Dialect} & \textbf{\vfive{} win} & \textbf{95\% CI}
  & \textbf{per game} & \textbf{\vfive{}} & \textbf{\vfour{}} \\
\midrule
default (\S\ref{sec:rules}) & \numDialectDefault\%    & \numDialectDefaultCI\%    & \numDialectDefaultForcedRate{}    & \numDialectDefaultForcedFive\%    & \numDialectDefaultForcedFour\% \\
no out-of-turn declaration & \numDialectNoOOT\%      & \numDialectNoOOTCI\%      & \numDialectNoOOTForcedRate{}      & \numDialectNoOOTForcedFive\%      & \numDialectNoOOTForcedFour\% \\
no cardless declaration & \numDialectNoCardless\% & \numDialectNoCardlessCI\% & \numDialectNoCardlessForcedRate{} & \numDialectNoCardlessForcedFive\% & \numDialectNoCardlessForcedFour\% \\
\vthree{} legacy & \numDialectLegacy\%     & \numDialectLegacyCI\%     & \numDialectLegacyForcedRate{}     & \numDialectLegacyForcedFive\%     & \numDialectLegacyForcedFour\% \\
\bottomrule
\end{tabular}
\end{table}

The head-to-head margin is insensitive to the dialect: the four rows of
Table~\ref{tab:dialects-v05} span \numDialectDefault\% to \numDialectNoOOT\%,
a range narrower than the width of any one of their intervals, and every interval
overlaps every other. None of them separates from $50\%$; the closest is the
no-out-of-turn row, whose interval reaches down to exactly
$50\%$. The v0.4 study's dialect experiment moved by
several points, but it measured a different pairing on a different bank, so the
two are not comparable; all that is claimed here is that the \vfive{}-against-
\vfour{} margin does not depend on which of these four dialects is played.

What the dialect does change is the population of states the mechanisms act on. Disabling out-of-turn declarations raises the
forced-endgame rate by a factor of about \numDialectForcedRatio{}, from
\numDialectDefaultForcedRate{} declarations per game to
\numDialectNoOOTForcedRate{}, because a team that may only declare on its own
turn arrives at the cardless endgame far more often. On that population
\vfive{}'s forced declarations are correct \numDialectNoOOTForcedFive\% of the
time against \vfour{}'s \numDialectNoOOTForcedFour\%, and \vfour{}'s overall
declaration accuracy falls to \numDialectNoOOTDeclFour\% against \vfive{}'s
\numDialectNoOOTDeclFive\%. The absolute accuracies are much higher than the
\numVfiveForcedAcc\% of Table~\ref{tab:forced} because the states are not the
same states --- a forced endgame reached because declaration was illegal until
now is far more determined than one reached after a full game of asking --- so
the two measurements are not comparable in level. Nor are they comparable in the
size of the gap between the two policies: the gap here, \numDialectNoOOTForcedFive\%
against \numDialectNoOOTForcedFour\%, is less than half the one
Table~\ref{tab:forced} records between \numVfiveForcedAcc\% and
\numVfourForcedAcc\%. What transfers across them is only the sign of that gap,
which is what M2 predicts, and that sign is the only claim made from this row. Ask hit rate also falls sharply under
these dialects, to \numDialectNoOOTHit\% from \numDialectDefaultHit\%, which is a
property of the dialect and not of either policy.

\subsection{Throughput}
\label{sec:results-throughput}

\vfive{} plays \numThroughputFive{} self-play games per second over a timing run
of \numThroughputGamesFive{} games, a wall-clock rate for one run with no
interval. The more useful comparison is the matched one the ablation bank
supplies, because it holds the opponent, the deals and the machine fixed: on that
bank the mechanism-disabled control of Table~\ref{tab:ablations-v05} runs at
\numThroughputPatho{} games per second and the shipped configuration at
\numThroughputFixed{}, a factor of \numThroughputRatioFive$\times$. (The two
figures differ from the self-play number above because they are played against
\vfour{} rather than a mirror.) The failure mode of \S\ref{sec:diag-baseline} is
therefore a compute cost as well as a play defect: the deadlock is played out move
by move and every event in it runs a belief update. The factor exceeds the event
ratio --- \numAblControlEvents{} events per game against \numAblFullEvents{} ---
which is consistent with the certificate accumulation \S\ref{sec:limitations}
notes, but this timing is not instrumented finely enough to isolate that. Every
evaluation in this section completes in minutes at these rates; the fit of
\S\ref{sec:fitting} does not, and its cost is not reported here.

\renewcommand{\topfraction}{0.7}
\renewcommand{\bottomfraction}{0.3}
\renewcommand{\textfraction}{0.2}
\renewcommand{\floatpagefraction}{0.5}
\setcounter{topnumber}{2}
\setcounter{bottomnumber}{1}
\setcounter{totalnumber}{3}
\setlength{\textfloatsep}{20pt plus 2pt minus 4pt}
\setlength{\floatsep}{12pt plus 2pt minus 2pt}
\setlength{\intextsep}{12pt plus 2pt minus 2pt}
